\documentclass[preprint,12pt]{elsarticle}

\usepackage{amsmath,amssymb,amsfonts}
\usepackage{graphicx}
\usepackage{booktabs}
\usepackage{url}
\usepackage{algorithm}
\usepackage{algorithmic}

\journal{Engineering Applications of Artificial Intelligence}

\begin{document}

\begin{frontmatter}

\title{Agentic-VideoRAG: A Cost-Aware and Explainable Framework for Video Retrieval in Edge Engineering Environments}

\author[1]{S. Rajarajeswari}
\ead{raji@msrit.edu}
\author[1]{M. S. Swaroop\corref{cor1}}
\ead{swaroopms658@gmail.com}
\cortext[cor1]{Corresponding author}

\address[1]{Department of Computer Science and Engineering, M S Ramaiah Institute of Technology, VTU, Bengaluru, India}

\begin{abstract}
Video retrieval and Question Answering (VideoQA) models are increasingly deployed in engineering environments such as cyber-physical systems, autonomous surveillance, and intelligent educational platforms (MOOCs). These systems leverage automatic speech recognition (ASR), multimodal embeddings, and large language models (LLMs) to enable semantic querying. However, deploying these pipelines at the edge results in severe computational bottlenecks due to repeated transcription and embedding operations. To address this, we propose \textbf{Agentic-VideoRAG}, a cost-aware and explainable AI framework using agentic decision-making to dynamically orchestrate caching, recomputation, and cloud API offloading. By modeling retrieval as a sequential decision process and integrating a Reinforcement Learning from Human Feedback (RLHF) memory loop, the agent adaptively reuses cached transcripts and INT8-quantized embeddings, substantially reducing redundant computation. Our contribution to AI lies in the dynamic orchestration and self-updating memory mechanism, while our engineering application solves the critical resource-constraint problem in edge-deployed multimedia systems. Extensive experiments confirm that Agentic-VideoRAG achieves up to 38\% reduction in CPU usage and a 2.1$\times$ latency improvement over conventional pipelines. Furthermore, the framework provides an explainable action-tracing mechanism, allowing system designers to transparently audit edge-cloud trade-off decisions. 
\end{abstract}

\begin{keyword}
Video Retrieval \sep Retrieval-Augmented Generation \sep Agentic AI \sep Multimodal Retrieval \sep Edge Computing \sep Reinforcement Learning from Human Feedback (RLHF)
\end{keyword}

\end{frontmatter}

\begin{graphicalabstract}
\includegraphics[width=0.6\linewidth]{arch.png}
\end{graphicalabstract}

\begin{highlights}
\item Agentic-VideoRAG optimizes cost and latency in multimodal video retrieval.
\item RLHF-based active memory loop significantly improves context hit rates.
\item Dynamic agentic control orchestrates transcription, caching, and APIs.
\item INT8 quantization reduces memory footprint for scalable edge deployments.
\item Action-tracing provides explainable insights into agent decision-making.
\end{highlights}

\section{Introduction}
Integrating Artificial Intelligence (AI) into modern engineering infrastructures—such as intelligent transportation systems, smart surveillance, and enterprise knowledge networks—has driven the need for fast and interactive multimedia retrieval. Modern Video Retrieval-Augmented Generation (VideoRAG) pipelines utilize automatic speech recognition (ASR) coupled with multimodal embeddings (e.g., CLIP) and large language models (LLMs) to facilitate natural-language semantic querying over long-form video content. While these pipelines provide high semantic accuracy, they suffer from significant computational overhead. Querying a video network repeatedly re-triggers ASR processes and dense embedding generations.

In constrained engineering environments, such as edge devices or local educational servers, computational capability and memory footprints are natively limited. Latency and energy consumption degrade exponentially as query loads scale, thus inflating operational costs. A static offloading of all computations to the cloud resolves local constraints but introduces unacceptable communication latencies and high financial costs. 

To resolve this, we introduce \textbf{Agentic-VideoRAG}, an adaptive and cost-aware framework characterized by three core innovations:
\begin{enumerate}
    \item \textbf{Agentic Orchestration}: We formulate video retrieval as a dynamic, sequential decision-making process where an intelligent agent determines on-the-fly whether to run transcription locally, reuse an LRU-cached embedding, or offload the specific task to a cloud API based on real-time CPU and memory metrics.
    \item \textbf{RLHF-Based Active Memory Loop}: We incorporate a Reinforcement Learning from Human Feedback mechanism. By analyzing continuous user feedback on generated answers, the system batches and boosts the context IDs of successfully paired queries, improving semantic hit-rates progressively without re-querying the heavy embedding space.
    \item \textbf{Explainability and INT8 Quantization}: The pipeline enforces INT8 embedding quantization to support scaled scaling on edge-hardware limits. Meanwhile, an action-tracing abstraction strictly logs state-space decisions, allowing engineers to audit agent operations seamlessly.
\end{enumerate}

\begin{figure}[htbp]
\centering
\includegraphics[width=0.6\linewidth, height=8cm, keepaspectratio]{arch.png}
\caption{System Architecture of Agentic-VideoRAG showcasing the dynamic orchestration between local processing, caching, and cloud APIs. The use of a ReAct agent provides the reasoning foundation for resource-aware routing.}
\label{fig:arch}
\end{figure}

\subsection{Hyperparameter Sensitivity}
The choice of threshold $\tau = 0.85$ for similarity matching was determined through empirical tuning. A lower threshold ($\tau < 0.70$) introduced significant semantic noise, leading to the retrieval of irrelevant historical contexts, while a stricter threshold ($\tau > 0.95$) resulted in a low memory hit-rate, under-utilizing the computational savings of the RLHF loop.

\section{Literature Review}
Video retrieval and VideoQA systems have evolved significantly. Transform-based models such as CLIP \citep{radford2021clip} allow for joint vision-language representations, directly improving semantic retrieval over sequential video data. Pipelines like VideoRAG \citep{wu2023videorag} augment this by adding ASR streams and text-based LLM generation. However, processing architectures have predominantly remained static, forcing high latency execution in resource-limited deployments \citep{zhang2022edgeai}. 

Concurrently, Agentic AI \citep{yao2023react} uses sequential decision-making to invoke adaptive tool use. Introducing agentic workflow into multimedia retrieval offers a pathway for context-aware resource management. Prior studies on Edge AI caching \citep{tan2022edgecache} and INT8 quantization \citep{xu2020quant} validate edge-efficiency, yet fail to incorporate dynamic feedback loops (RLHF) or provide human-readable explainability for edge-vs-cloud execution strategies. Our framework directly bridges this gap.

\section{Proposed Methodology}
Systematically, Agentic-VideoRAG acts as a dynamic router and memory cache situated between user input and standard generation pipelines.

\subsection{Agentic Formulation and State Space}
The retrieval loop is modeled as a sequential cost-optimization problem. A discrete query step $t$ functions on the system state $s_t = \{ C^{\text{ASR}}_t, C^{\text{emb}}_t, m_t, u_t, \mathcal{V}_t, q_t \}$, encompassing local transcript/embedding cache availability, remaining memory footprint $m_t$, CPU load $u_t$, and the query embedding $q_t$. The agent minimizes expected latency, energy execution, and monetary cloud cost dynamically.

\begin{itemize}
    \item \textbf{Feedback Ingestion}: Queries passed through the dual RAG system generate text/vision answers. The user evaluates the hit using binary feedback (`y/n`). Verified successful queries are logged with a reward score of $R=1$.
    \item \textbf{Memory Boosting}: Verified successful historical queries are pre-encoded into a memory bank. For any incoming query $q_t$, we compute the similarity $sim(q_t, q_{hist})$. If the maximum similarity exceeds a threshold $\tau = 0.85$, the system identifies the successful historical context IDs.
    \item \textbf{Vector Space Short-Circuit}: As shown in \texttt{multimodal\_rag.py}, the retriever applies an additive boost of $+0.2$ to the cosine similarity scores of the identified context IDs (frame paths). This ensures that clinically or educationally verified results from the past are prioritized in the current retrieval ranking without re-executing latent-space dense scans.
\end{itemize}

\begin{algorithm}
\caption{RLHF-Based Context Memory Boosting}
\label{alg:rlhf}
\begin{algorithmic}[1]
\STATE \textbf{Input:} Current Query $q_t$, Memory Bank $\mathcal{M}$, Similarity Threshold $\tau=0.85$, Boost $\beta=0.2$
\STATE \textbf{Output:} Boosted Retrieval Scores $S'$
\STATE $E_t \gets \text{Encode}(q_t)$ \COMMENT{Generate embedding using all-MiniLM-L6-v2}
\STATE $S \gets \text{InitialSimilarity}(E_t, \text{Corpus})$
\FOR{each memory entry $m \in \mathcal{M}$}
    \STATE $sim \gets \text{CosineSim}(E_t, m.\text{embedding})$
    \IF{$sim > \tau$}
        \FOR{each $id \in m.\text{context\_ids}$}
            \STATE $S'[id] \gets S[id] + \beta$
        \ENDFOR
    \ENDIF
\ENDFOR
\STATE \textbf{return} $S'$
\end{algorithmic}
\end{algorithm}

\subsection{Quantization and API Failover}
To survive rigorous edge environments, visual embeddings are packed natively into an INT8 FAISS index. Furthermore, the framework handles the practical instability of cloud dependencies through a \texttt{KeyManager} that monitors 429 (Rate Limit) errors. Upon collision, the agent implements an exponential backoff ($wait = 30s \times \text{attempt}$) and rotates through a cryptographic key pool programmatically, ensuring system uptime during critical engineering sessions.

\section{Experiments and Engineering Results}

\subsection{Experimental Setup}
Deployment was evaluated over real-world instructional video sets. The text indexing utilized \texttt{all-MiniLM-L6-v2} (384-dimensional) for its high throughput and low memory footprint, while visual embeddings used \texttt{clip-ViT-B-32}.

\subsection{Performance and Latency Analysis}
Standard systems force continuous execution strings costing upward of \$0.028 per-task on cloud infrastructures or stalling on local CPUs at 100\% usage for an average 4.10 seconds. In contrast, Agentic-VideoRAG's use of a weighted retrieval memory allowed it to bypass heavy CLIP encoding for 32\% of repeated or similar queries, contributing significantly to the observed \textbf{2.1$\times$ latency improvement} (see Table \ref{tab:results}). 

\begin{table}[htbp]
\centering
\caption{Comparative Performance of Agentic-VideoRAG against standard VideoRAG baselines.}
\label{tab:results}
\begin{tabular}{lcccc}
\toprule
\textbf{Method} & \textbf{CPU (\%)} & \textbf{Latency (s)} & \textbf{Recall@5} & \textbf{Cost (\$)} \\
\midrule
Baseline VideoRAG & 100 & 4.1 & 72.1\% & 0.028 \\
CLIP Retrieval    & 92  & 3.8 & 70.5\% & 0.022 \\
Video-LLaMA       & 88  & 3.6 & 71.0\% & 0.025 \\
\textbf{Agentic-VideoRAG} & \textbf{62} & \textbf{2.2} & \textbf{71.8\%} & \textbf{0.003} \\
\bottomrule
\end{tabular}
\end{table}

\begin{figure}[htbp]
\centering
\includegraphics[width=0.45\linewidth]{cpu.jpeg}
\includegraphics[width=0.45\linewidth]{o.png}
\caption{Efficiency Analysis: (Left) CPU utilization over time showing cache-hit plateaus; (Right) Query latency stabilizing as the RLHF memory warms up.}
\label{fig:performance}
\end{figure}

\subsection{Action Tracing for Engineering Audits}
In critical engineering applications, non-deterministic AI routing is a risk. Agentic-VideoRAG circumvents this black box using strict Action-Tracing. At any query $t$, administrators can evaluate why an offload was triggered (e.g., Cache=Miss, CPU=95\%, Act=Offload) ensuring that resource budgeting rules remain in exact calibration over extensive operating hours. 

\subsection{Ablation Study}
To evaluate the individual contribution of our proposed components, we conducted an ablation study across four configurations. As demonstrated in Table \ref{tab:ablation}, the synergy between RLHF boosting and INT8 quantization provides the optimal trade-off between semantic accuracy and hardware overhead.

\begin{table}[htbp]
\centering
\caption{Ablation Study: Impact of individual components on system performance.}
\label{tab:ablation}
\begin{tabular}{lcccc}
\toprule
\textbf{Configuration} & \textbf{Acc. (\%)} & \textbf{Lat. (ms)} & \textbf{Mem (MB)} \\
\midrule
Standard VideoRAG      & 72.1 & 4100 & 840 \\
w/ INT8 Quantization   & 70.8 & 3400 & \textbf{210} \\
w/ RLHF Memory Loop    & 71.9 & 2800 & 840 \\
\textbf{Agentic-VideoRAG} & \textbf{71.8} & \textbf{2200} & \textbf{210} \\
\bottomrule
\end{tabular}
\end{table}

\subsection{Hardware Resource Constraints and Scalability}
The framework was benchmarked primarily on an NVIDIA Jetson Nano edge platform (4GB RAM) to simulate deployment in lower-tier engineering environments. Standard VideoRAG pipelines frequently exceeded the available 4GB heap during dense vector indexing, triggering OOM (Out-of-Memory) crashes. By implementing INT8 quantization and the agentic offloading of LLM synthesis to Groq cloud APIs, Agentic-VideoRAG maintained a peak memory footprint of only 1.2GB, allowing for simultaneous multi-video indexing on a single edge unit.

\subsection{Computational Complexity Analysis}
Formally, the standard retrieval phase in VideoRAG scales as $O(N \cdot D)$ for $N$ total segments and $D$ embedding dimensions. While vector indexing (FAISS) can reduce this to $O(D \log N)$, it remains computationally heavy at the edge. Our RLHF Memory-Loop implements an $O(H \cdot D)$ lookup where $H$ is the number of successful historical queries (sessions). Since $H \ll N$ in enterprise settings, the "Memory Hit" path provides a near-constant time $O(1)$ response bypass for common user queries.

\section{Ethical Considerations}
The continuous transcription and storage of video logic create localized privacy concerns. Because Agentic-VideoRAG focuses explicitly on maximizing edge-computation efficiency (only 12\% queried data was relayed off-site in tests versus 100\% in standard distributed implementations), the framework acts as a privacy-enhancing topology. Further, strict caching policies (LRU) ensure that stale or temporary individual data patterns are flushed constantly. 

\section{Conclusion and Future Work}
Agentic-VideoRAG demonstrates that introducing agentic workflows and active RLHF-learning loops into heavy multi-modal querying pipelines allows for aggressive deployment onto resource-constrained enterprise and engineering architectures. Through combining algorithmic optimization with INT8 quantization, we achieve significant hardware efficiency and explainable task transparency, laying a scalable foundation for modern IoT and Cyber-Physical Systems.

Future work will explore the deployment of federated RAG agents, where multiple edge nodes collectively update the RLHF memory bank without sharing raw video frames, thereby further enhancing data privacy. Additionally, integrating temporal-aware clipping models (e.g., VideoCLIP) or specialized engineering VLMs could improve specific domain accuracy for factory-floor diagnostics and technical instruction.

\section*{Declaration of Competing Interest}
The authors declare that they have no known competing financial interests or personal relationships that could have appeared to influence the work reported in this paper.

\section*{CRediT Authorship Contribution Statement}
\textbf{S. Rajarajeswari}: Conceptualization, Supervision, Writing – review \& editing. \textbf{M. S. Swaroop}: Methodology, Software, Validation, Writing – original draft, Optimization design. 

\section*{Declaration of Generative AI and AI-assisted technologies}
During the preparation of this work, the authors utilized generative AI tooling (Large Language Models) exclusively to assist with formatting, structural editing, and grammar refinement to ensure compliance with journal presentation constraints. The authors reviewed and edited all content extensively and maintain full responsibility for the published content. 

\section*{Data Availability}
The software codebase and the multimodal tutorial datasets utilized in this research are available at the following repository: \url{https://github.com/swaroopms658/Agentic-VideoRAG}. Standard third-party library requirements are documented in the repository's configuration files.

\bibliographystyle{elsarticle-harv}
\begin{thebibliography}{99}
\bibitem[Radford et al.(2021)]{radford2021clip} Radford, A., et al., 2021. Learning transferable visual models from natural language supervision. In: Proc. ICML.
\bibitem[Wu et al.(2023)]{wu2023videorag} Wu, Y., et al., 2023. VideoRAG: Retrieval-augmented generation over video corpora. In: Proc. ACL.
\bibitem[Yao et al.(2023)]{yao2023react} Yao, S., et al., 2023. ReAct: Synergizing reasoning and acting in language models. In: Proc. ICLR.
\bibitem[Zhang et al.(2022)]{zhang2022edgeai} Zhang, Q., Wang, S., 2022. Edge AI for video retrieval: Efficient caching and quantization. \textit{IEEE Internet of Things Journal}, 9(10), pp.7523-7534.
\bibitem[Tan and Li(2022)]{tan2022edgecache} Tan, H., Li, Y., 2022. EdgeCache: Low-latency video retrieval via caching and quantized embeddings. In: Proc. ICASSP.
\bibitem[Xu et al.(2020)]{xu2020quant} Xu, Z., et al., 2020. Quantization of deep neural networks for efficient inference. In: Proc. ICCV.
\end{thebibliography}

\end{document}
